\section{Related Work}
\label{sec:related}

\para{Model Collapse.} The foundational concept of model collapse was formalized by \citet{shumailov2024ai}, who demonstrated that recursively training models on synthetic data leads to a severe loss of the original data distribution's tails. Their work primarily evaluates single-model recursive loops, showing a baseline ``collapse'' scenario when a population effectively has a size of one. Building on this, \citet{schaeffer2025model} provided a meta-analysis critiquing the inevitability of model collapse, suggesting that in real-world scenarios where models are fine-tuned on a mix of human and AI data, collapse might be highly avoidable. Unlike these single-agent or mixed-data scenarios, we focus strictly on multi-agent synthetic ecosystems to isolate the effect of population size and diversity.

\para{Multi-Agent Ecosystems and Cooperation.} Recent research has extended the study of model collapse to multi-agent environments. \citet{piatti2024cooperate} look at collapse from a resource management perspective, simulating societies of LLM agents managing a common pool resource. They found that without advanced reasoning, most LLM societies exhaust resources and collapse, emphasizing the need for diversity and communication. While their work focuses on agent interactions, we focus on the data-driven model collapse resulting from shared synthetic training pools.

\para{Epistemic Diversity.} Directly addressing the minimum viable population hypothesis, \citet{hodel2025epistemic} evaluated multiple distinct LLMs forming an ecosystem. They found that increasing AI ecosystem diversity mitigates knowledge collapse, as models with different ``epistemic'' viewpoints act as a balancing force. We complement this work by providing an empirical comparison of homogeneous versus heterogeneous populations under the same recursive training framework, quantifying how even simple diversity mechanisms like varying temperatures can serve as a counter-weight to collective bias.
